\chapter{Ce que voyait Kamatsu}
\label{chap:ce-que-voyait-kamatsu}

Kamatsu avait grandi sur les routes.

À huit ans, la caravane était pour lui à la fois une maison et un immense
terrain de jeu. Il connaissait les voyageurs qui la composaient, les chevaux,
les chiens, les chariots et jusqu'aux habitudes de certains animaux. Lorsque
la route était tranquille et qu'on l'autorisait à circuler, il était rare
qu'il reste longtemps auprès de ses parents.

Il passait d'un chariot à l'autre, marchait quelque temps auprès de l'un,
grimpait sur un autre lorsqu'on l'y invitait, puis repartait dès que quelque
chose attirait son attention. Il discutait avec les adultes, jouait avec les
autres enfants et s'arrêtait parfois pour caresser un chien avant de courir
rattraper le reste de la caravane.

Ce jour-là, pourtant, ce furent les bruits du voyage qui occupèrent son esprit.

Les sabots des chevaux frappaient régulièrement le sol. Les roues grinçaient
sur leurs essieux et cahotaient chaque fois qu'elles rencontraient une pierre.
Les harnais tintaient au rythme des pas, accompagnés par les conversations des
voyageurs et les appels que l'on lançait parfois d'un chariot à l'autre.

Kamatsu écouta.

Après quelques instants, il commença à entendre autre chose.

Les sabots revenaient toujours presque au même moment. Le grincement d'une roue
se glissait entre deux séries de pas. Un cheval soufflait régulièrement un peu
plus loin et, chaque fois que la route devenait mauvaise, les caisses d'un
chariot s'entrechoquaient.

Kamatsu posa une main contre le flanc de bois du chariot auprès duquel il
marchait.

Il frappa deux fois.

Puis attendit.

Les chevaux continuèrent.

Kamatsu frappa de nouveau.

Deux coups rapides, un silence, puis un troisième.

Une voyageuse assise sur le chariot se pencha pour le regarder.

— Qu'est-ce que tu fais ?

— J'écoute.

Elle observa autour d'elle.

— Quoi ?

Kamatsu désigna successivement les chevaux, les roues et les harnais.

— Tout.

Il frappa encore le bois, cette fois juste après le claquement d'une roue
contre une pierre.

La femme sourit.

— Et tu es obligé de participer ?

Kamatsu réfléchit sérieusement à la question.

— Il manquait quelque chose.

Elle éclata de rire.

Kamatsu reprit aussitôt son étrange accompagnement. Il changeait parfois de
rythme, essayait de répondre au bruit des roues ou attendait plusieurs secondes
avant de recommencer. Cela ne ressemblait probablement à une mélodie que pour
lui.

Mais cela lui suffisait.

Il continua ainsi quelque temps, frappant le bois au rythme de la caravane, jusqu'à ce que son regard se perde dans le ciel.

Un nuage avait une forme étrange.

Il ralentit pour mieux l'observer. Avec un peu d'imagination, ce n'était plus
tout à fait un nuage. Il avait une tête, un dos immense et quelque chose qui
pouvait presque passer pour une queue.

Un dragon.

Kamatsu continua de marcher en levant les yeux vers le ciel. Le dragon avançait
lentement au-dessus de la caravane, poussé par le vent. Sa gueule s'allongeait,
une aile disparaissait déjà, mais il suffisait de regarder autrement pour la
retrouver.

Kamatsu ne vit pas la flaque. Son pied s'enfonça dedans jusqu'à la cheville, projetant une gerbe d'eau boueuse sur son pantalon. Il baissa enfin les yeux tandis que quelques rires éclataient autour de lui.

Kamatsu regarda sa jambe, puis le ciel. Le dragon avait déjà commencé à se
défaire.

— Il y avait un dragon.

— Bien sûr, répondit quelqu'un.

— Dans le ciel.

— Évidemment.

Kamatsu désigna le nuage, mais celui-ci ne ressemblait déjà plus vraiment à
rien.

Il fronça les sourcils.

— Il était là.

Les rires reprirent tandis qu'il courait rattraper le chariot de ses parents.

\scenebreak

Le soir, les chariots formaient le camp et les bruits de la route cédaient peu
à peu la place à ceux du feu.

Kamatsu avait mangé avec ses parents, mais n'était pas resté longtemps auprès
d'eux. Plusieurs petits groupes s'étaient formés autour des flammes et il
passait de l'un à l'autre avec la même facilité qu'il passait d'un chariot à
l'autre pendant la journée.

Il connaissait presque tout le monde.

Et presque tout le monde le connaissait.

Kamatsu s'assit d'abord auprès de quelques voyageurs qui discutaient de leur
prochaine étape. Il les écouta quelques instants avant de profiter du premier
silence.

— Tu connais une histoire ?

L'homme auquel il venait de s'adresser le regarda.

— Encore ?

— Une que tu ne m'as jamais racontée.

— Je crois que tu les connais toutes.

Kamatsu secoua immédiatement la tête.

— Non.

— Comment peux-tu le savoir ?

— Parce qu'il y en a toujours d'autres.

Quelques personnes sourirent autour du feu.

L'homme finit par céder et commença à raconter une vieille histoire qu'il
tenait d'un voyageur rencontré plusieurs années auparavant. Kamatsu l'écouta,
les yeux fixés sur lui, interrompant parfois le récit pour demander pourquoi
quelqu'un avait pris telle décision ou ce qui s'était passé ensuite.

Puis son regard revint vers les flammes.

Il resta silencieux quelques secondes.

— Tu le vois ?

Le conteur s'interrompit.

— Quoi ?

Kamatsu désigna le feu.

— Le cheval.

Plusieurs regards se tournèrent vers les flammes.

— Quel cheval ?

— Là.

Kamatsu pointa une langue de feu qui s'élevait au-dessus d'une bûche.

— Sa tête est là. Et là, c'est son dos.

La flamme vacilla.

— Maintenant, il court.

Un des voyageurs plissa les yeux.

— Moi, je vois surtout du feu.

— Il faut regarder mieux.

Cette réponse provoqua quelques rires.

La forme disparut presque aussitôt, remplacée par une autre lorsque la bûche
s'effondra légèrement sur elle-même.

Kamatsu pencha la tête.

— Maintenant, c'est plus un cheval.

— Ah ?

— Non.

Il réfléchit.

— On dirait une tour.

Quelqu'un ajouta une branche au feu et la tour disparut à son tour.

— Tu l'as cassée !

— Toutes mes excuses.

Kamatsu regarda encore quelques instants les flammes, attendant qu'une nouvelle
forme y apparaisse.

— Et après ?

Le voyageur sourit et reprit son histoire.

Kamatsu resta auprès du feu jusqu'à la fin. Puis il en rejoignit un autre, un
peu plus loin, où il trouva de nouvelles personnes à écouter et de nouvelles
histoires à réclamer. Il aurait probablement continué ainsi encore longtemps
si une voix familière ne l'avait pas rappelé à l'ordre.

— Kamatsu.

Il tourna la tête. Sa mère se tenait à quelques pas de lui.

— Il est temps d'aller te coucher.

— Maintenant ?

— Maintenant.

Kamatsu regarda le feu, puis les voyageurs qui l'entouraient, comme s'il
cherchait parmi eux un soutien à sa cause.

— Encore une histoire.

Sa mère eut un petit rire.

— Encore une ?

— Une dernière.

— C'est ce que tu as dit tout à l'heure.

— Mais celle-là, ce sera vraiment la dernière.

— Et demain ?

Kamatsu réfléchit un instant.

— Demain, ce seront celles de demain.

Quelques rires accompagnèrent sa réponse. Sa mère secoua la tête, amusée, puis
lui tendit la main.

— D'accord. Une dernière. Mais dans le chariot.

Kamatsu se leva aussitôt, avant qu'elle puisse changer d'avis.

\scenebreak

Sa mère le raccompagna jusqu'au chariot des Shinzawa. Kamatsu se glissa parmi
les couvertures tandis qu'elle s'installait près de lui.

— Alors ?

— Alors quoi ?

— L'histoire.

Sa mère sourit.

Elle commença à lui raconter l'histoire d'un voyageur qui avait quitté son
village pour découvrir ce qui se trouvait derrière les montagnes. Kamatsu
l'interrompit presque immédiatement pour demander pourquoi il était parti,
s'il voyageait seul et ce qu'il espérait trouver de l'autre côté.

Sa mère répondit à ses premières questions, puis lui demanda de la laisser
continuer s'il voulait connaître la suite.

Cette fois, Kamatsu se tut.

Il écouta le voyageur franchir les premières collines, traverser une rivière et
s'engager sur un chemin qu'il ne connaissait pas. Dehors, les dernières
conversations de la caravane s'éteignaient peu à peu, mêlées au crépitement
lointain des feux et aux bruits des animaux installés pour la nuit.

La voix de sa mère devint bientôt le seul son auquel Kamatsu prêtait encore
attention.

Il n'entendit jamais la fin de l'histoire.